% Target language: German (Germany), de-DE.
% Translation source: content/en-US.tex at repository commit 9e0108924aa495091c8ea8a768077439f3c31740.
% Checked against Simplified Chinese source commit 0d064f6ce5c6f02a2e89edc6004b56723e8042be.
% Status: AI-assisted translation; needs native-speaker review.
\StoryParagraph{3}{Das Gewicht einer einzelnen Eizelle: 3,5 Mikrogramm.}

\StoryParagraph{3}{Das Gewicht von fünfzig Millionen Dollar in bar: 2,5 Tonnen.}

\StoryParagraph{3}{Bei einem Anruf aus dem Montage Laguna Beach verlangte Jing Tian von mir Letzteres, wobei Ersteres als Sicherheit diente.}

\StoryParagraph{3}{Sie war dorthin gereist, um sich Eizellen entnehmen zu lassen. Das Kind sollte von einer Leihmutter ausgetragen werden und 2027 zur Welt kommen. Wenn alles schnell genug ging, konnte es noch im Jahr des Pferdes geboren werden. Sie selbst hatte das Thema angesprochen. In der Klinik sagte sie: keine fünfzig Millionen, keine Entnahme.}

\UnprintedPause

\StoryParagraph{0}{Das Montage Laguna Beach liegt auf einer Klippe über dem Pazifik. Wir mieteten die gesamte Etage für dreißig Tage, eine Million Dollar, um ihre Privatsphäre zu schützen. Nach acht Tagen reiste sie ab.}

\StoryParagraph{0}{Während der verbleibenden zweiundzwanzig Tage gingen die Zimmermädchen des Hotels weiterhin jeden Tag wie gewohnt hinauf. Meine Assistentin erzählte mir davon. Sie hatte es einmal gesehen, als sie den Schlüssel zurückbrachte: Eine Mexikanerin schob ihren Wagen von Zimmer zu Zimmer, wechselte die Handtücher, wischte die Ablagen ab, steckte die Bettlaken unter die Matratzen, trat wieder hinaus und zog jede Tür hinter sich zu. Zwei Stunden für die ganze Etage.}

\StoryParagraph{0}{Meine Assistentin fragte sie, warum sie das noch immer tat, obwohl niemand dort wohnte.}

\StoryParagraph{0}{Sie sagte, im Belegungsplan stehe, dass jemand dort wohne.}

\StoryParagraph{0}{Versailles, seiner Königin beraubt.}

\StoryParagraph{0}{Da begriff ich plötzlich, dass Jing Tian so etwas mochte. Jing Tian mochte es, wenn ein Ort ihretwegen leer stand. Wenn nichts geschah.}

\StoryParagraph{3}{Wie das Montage Laguna Beach stand auch eine Gulfstream G700 leer. Sie parkte verloren in Van Nuys, während ihre Besatzung jeden Tag wie gewohnt zum Dienst erschien.}

\StoryParagraph{3}{Am achten Tag führte sie jenes Telefonat. Danach sah sie niemand je wieder.}

\NarrativeBreak

\StoryParagraph{0}{Schon 2007 verbrannte jemand Geld für sie. In jenem Jahr hatte ich gerade mein Studium an der Peking-Universität begonnen.}

\StoryParagraph{3}{Natürlich war ich nicht derjenige, der das Geld verbrannte.}

\StoryParagraph{3}{Bei QQ öffnete sich ein Fenster mit ihr als Werbegesicht. Es war auf jenem alten Lenovo-IBM-Rechner, die Auflösung völlig verschwommen. Ich speicherte das Bild.}

\StoryParagraph{3}{In meinem Wohnheimzimmer an der Peking-Universität, das tausend Yuan im Jahr kostete, ging ich ins Netz, fand sie auf Xiaonei und schickte ihr eine Freundschaftsanfrage. Bei Xiaonei musste man eine Nachricht hinzufügen. Etwa vierzig Minuten lang schrieb und löschte ich, löschte und schrieb neu. Am Ende schickte ich: „Hallo.“}

\StoryParagraph{3}{Sie nahm die Anfrage nicht an.}

\StoryParagraph{3}{Vielleicht hatte sie Xiaonei nie benutzt.}

\StoryParagraph{3}{In den Jahren danach dachte ich bei allem, was ich tat, über eine Frage nach: Wie weit musste ich es bringen, damit sie mich annahm? Bei meiner ersten Milliarde dachte ich daran. Bei zehn Milliarden dachte ich wieder daran.}

\StoryParagraph{3}{Irgendwann hörte ich auf, darüber nachzudenken. Renren war längst pleitegegangen.}

\StoryParagraph{3}{Eine zweite Anfrage schickte ich nie.}

\StoryParagraph{3}{2011 war es kalt in Peking. Ich kaufte eine Daunenjacke für 150 Yuan. Ich erinnere mich an den Preis, weil ich drei Tage lang gezögert hatte.}

\StoryParagraph{3}{Im selben Jahr spielte sie die Hauptrolle in dem 150 Millionen Yuan teuren Film \textit{The Warring States}. Sun Honglei, Francis Ng, Kim Hee-sun, Kiichi Nakai und Jiang Wu waren als bloße Staffage herbeigeholt worden, genau wie neunzehn Jahre später das Montage Laguna Beach und die Gulfstream G700. Sie wussten nicht, warum sie dort waren. Sie waren so weit gekommen, also musste doch etwas geschehen. Nichts geschah.}

\StoryParagraph{3}{Dann kam \textit{The Great Wall} mit einem Budget von 150 Millionen Dollar. Zhang Yimou. Andy Lau. Matt Damon. Dann Jackie Chan. Die Namensliste des Jing-Tian-Universums wurde immer länger.}

\StoryParagraph{3}{Damals lebte ich bereits in den Vereinigten Staaten. Ich studierte, gründete ein Unternehmen und hatte immer mehr zu tun. Nur eines änderte sich nie: Bei jedem Gerätewechsel übertrug ich dieses QQ-Bild, wie umständlich es auch war.}

\StoryParagraph{3}{Neunzehn Jahre später kam mein Name auf die Liste.}

\NarrativeBreak

\StoryParagraph{3}{Wir trafen uns zum ersten Mal im Rosewood Hong Kong. Ich hatte immer gedacht, ich würde an dem Tag, an dem ich sie traf, nervös sein.}

\StoryParagraph{3}{Das war ich nicht.}

\StoryParagraph{3}{Sie war dünner als auf den Fotos, lässig gekleidet und saß zurückgelehnt auf einem Sofa. Ich sah sie an und konnte mich plötzlich nicht mehr erinnern, worauf genau ich so lange gewartet hatte.}

\StoryParagraph{3}{Ich nahm sie mit zu Cattelans Banane. Sie betrachtete sie lange und fragte dann: „Wie viel kostet die?“}

\StoryParagraph{3}{6,2 Millionen Dollar.}

\StoryParagraph{3}{Sie sagte: „Was passiert, wenn sie verfault?“}

\StoryParagraph{3}{Ich sagte nichts. Sie sagte: „Lass gut sein. Gehen wir ins Kino.“}

\StoryParagraph{3}{Ich sagte: „Gut.“}

\StoryParagraph{3}{Sie sagte, im Kinosaal dürfe außer uns beiden niemand sein.}

\StoryParagraph{3}{Ich sagte: „Gut.“}

\StoryParagraph{3}{Ich ging hinaus, um zu telefonieren. Als ich zurückkam, stand sie noch immer vor der Banane.}

\StoryParagraph{3}{Der Film war \textit{Zoomania 2}.}

\StoryParagraph{3}{Ein Kinosaal für zwei. So etwas gibt es sonst gewöhnlich nur, wenn die Zentrale Propagandaabteilung einen Film prüft.}

\StoryParagraph{3}{Ich ließ meine Assistentin jeden Sitzplatz im Kino kaufen.}

\StoryParagraph{3}{Ich hatte Artikel gelesen, die manipulierte Einspielergebnisse anprangerten: alle Karten verkauft, aber nur zwei Menschen im Saal. Nun verstand ich. Vielleicht ging es nicht immer um die Einspielergebnisse.}

\StoryParagraph{3}{Alle Plätze zu kaufen war nicht schwer. Das Problem waren die bereits verkauften Karten. Meine Assistenten und Leibwächter kamen früh und warteten am Eingang. Wann immer jemand mit einer Karte auftauchte, drückten sie ihm Bargeld in die Hand.}

\StoryParagraph{3}{Ein Paar kam mit seinem Kind. Das Kind stand schon am Eingang und hielt einen Eimer Popcorn hoch in die Luft.}

\StoryParagraph{3}{Eine Assistentin hielt ihnen ein Bündel Geldscheine hin. Das Kind verstand nicht, warum sie nicht hineindurften. Die Eltern nahmen das Geld, bedankten sich, fragten nichts und gingen.}

\StoryParagraph{3}{An diesem Tag nahmen sechsundzwanzig Menschen das Geld.}

\StoryParagraph{3}{Nach dem Film ersetzte sie ihr WeChat-Profilbild durch Officer Judy Hopps.}

\StoryParagraph{3}{An diesem Abend fasste ich ihr an die Brüste. Sie stieß mich weg.}

\StoryParagraph{3}{Ich dachte, sie wolle das nicht. Ich zog die Hand zurück und wollte gerade etwas sagen.}

\StoryParagraph{3}{Sie sagte nichts. Sie richtete sich auf, holte aus ihrer Tasche einen kleinen Beutel und öffnete den Reißverschluss. Darin lagen in einer Reihe kleine Dinge, die zusammen wie ein winziger Werkzeugsatz aussahen.}

\StoryParagraph{3}{Sie sagte: „Deine Fingernägel tun mir weh.“}

\StoryParagraph{3}{Ich sah sie mir an. Ich hatte sie immer irgendwie mit dem Nagelknipser abgeschnitten. Die Kanten waren scharf. Normalerweise bemerkte ich es nicht.}

\StoryParagraph{3}{Sie sagte: „Komm her.“}

\StoryParagraph{3}{Sie legte meine Hand auf ihren Schoß und feilte einen Nagel nach dem anderen.}

\StoryParagraph{3}{Zuerst mit der groben Seite. Als sie damit fertig war, wechselte sie die Seite und dann noch einmal. Zwischendurch hielt sie einen Finger ins Licht, betrachtete ihn und feilte weiter.}

\StoryParagraph{3}{Ich fragte: „Wie lange dauert das?“}

\StoryParagraph{3}{Sie sagte: „Bis sie niemanden mehr kratzen.“}

\StoryParagraph{3}{Zehn Finger. Sie brauchte fast eine Stunde.}

\StoryParagraph{3}{Währenddessen erklärte sie mir: Diese Seite bringt sie in Form, diese poliert sie, und mit der letzten Seite darfst du nicht hin und her reiben. Nur in eine Richtung.}

\StoryParagraph{3}{Ich sagte: „Gut.“}

\StoryParagraph{3}{Als sie fertig war, drehte sie meine Hand um, betrachtete sie und strich mit den Fingerspitzen darüber.}

\StoryParagraph{3}{Sie sagte: „So.“}

\StoryParagraph{3}{„Nur zu. Fass an, wo du willst.“}

\StoryParagraph{3}{Zwischendurch sagte sie einen Satz. Sie sagte: „Wenn ich eines Tages nicht mehr da bin, wird dir nie wieder jemand die Fingernägel polieren.“}

\StoryParagraph{3}{Ich lächelte. Ich sagte nichts.}

\StoryParagraph{0}{Drei Monate nachdem ich sie kennengelernt hatte, ließ ich ein Dossier zusammenstellen. Ihre Projekte aus all den Jahren, Geschäftspartner, Unternehmensbeteiligungen, Jing Guoqing, Tian Xin’ai, alle zugehörigen Konten, die gemietete Wohnung in Peking, die gekaufte Wohnung in Shanghai und einige andere Dinge. Es dauerte eine Woche. Insgesamt etwas mehr als vierzig Seiten. Auch Geschichten über Männer und Frauen standen darin. Die hielt ich für das Unwichtigste.}

\StoryParagraph{0}{Einiges daraus erzählte sie mir später selbst. Ihre Version wich ein wenig von dem ab, was dort stand. Ich sagte nichts.}

\StoryParagraph{0}{Manches erzählte sie mir nie. Ich fragte auch nicht danach.}

\StoryParagraph{0}{Ich glaubte alles, was sie sagte.}

\StoryParagraph{0}{Ich bewahrte das Dossier auf. Ich habe es noch heute.}

\StoryParagraph{3}{Aber ich glaubte kein einziges Wort.}

\NarrativeBreak

\StoryParagraph{3}{Sie sagte: „In Hongkong sind zu viele Menschen. Jemand wird mich erkennen. Lass uns irgendwohin fahren.“}

\StoryParagraph{3}{Ich sagte: „Gut.“}

\StoryParagraph{3}{Sie tippte auf der Karte Teneriffa an.}

\StoryParagraph{3}{Ich sagte: „Gut.“}

\StoryParagraph{3}{Ich ging hinaus, um zu telefonieren. Erst als ich zurückkam, erfuhr ich, dass der Ort im Atlantik vor der Nordwestküste Afrikas lag. Am nächsten Tag um zwölf Uhr brach unsere dreißigköpfige Gruppe auf. Am Flughafen wartete ein A330.}

\StoryParagraph{3}{Im Schlafzimmer des Flugzeugs fragte sie mich: „Klingt mein Name schön?“}

\StoryParagraph{3}{Ich sagte: „Ja.“}

\StoryParagraph{3}{Sie sagte: „Er setzt sich aus dem Familiennamen meines Vaters und dem meiner Mutter zusammen. Ist das nicht ein erstaunlicher Zufall?“}

\StoryParagraph{3}{Ich sagte: „Ich habe ihn mir schon am ersten Tag gemerkt.“}

\StoryParagraph{3}{Sie legte die Arme um mich und sagte: „Nenn mich nicht mehr Jing Tian. Nenn mich Mama.“}

\StoryParagraph{3}{„Gut, Mama.“}

\StoryParagraph{3}{Nach der Landung kam der Kapitän zu mir. „Chef, die Insel wollte uns zuerst keine Landeerlaubnis geben. Sie hatten noch nie ein so großes Privatflugzeug gesehen.“}

\StoryParagraph{3}{Auf der Insel galten dieselben Regeln: Niemand durfte in unserer Nähe sein. Meine Assistenten kümmerten sich darum. Niemand fragte nach dem Grund. Diesmal fragte auch ich nicht, wie viele es waren.}

\StoryParagraph{3}{Das Ritz-Carlton blickte auf den Atlantik; seine bodentiefen Fenster waren zu groß. Sie konnte bei der kleinsten Spur Licht nicht schlafen. Die Assistenten klebten das ganze Zimmer mit schwarzem Klebeband ab. Draußen war es taghell, drinnen so dunkel wie in einer Isolierzelle. In der Sonne löste sich das Klebeband und musste jeden Tag neu angebracht werden.}

\StoryParagraph{3}{Es waren insgesamt elf Assistenten. Jeden Morgen kamen sie um 6.30 Uhr und gingen erst, nachdem wir eingeschlafen waren. Zwei junge Frauen waren für das Klebeband zuständig: Eine klebte die Fensterrahmen ab, die andere die Leuchtanzeigen der Klimaanlage und der Steckdosen. Wenn sie fertig waren, kontrollierten sie alles dreimal. Sie kauften dreißig Rollen Klebeband, später noch einmal zwanzig. Sie fanden heraus, dass Klebeband, das nach drei Uhr nachmittags angebracht wurde, sich am leichtesten löste, weil das Glas dann am heißesten war. Also brachten sie es fortan morgens an. Niemand hatte ihnen das beigebracht. Sie fanden es selbst heraus. Sie berieten sich darüber, als ginge es um ein Handwerk.}

\StoryParagraph{3}{Sie kannte die Namen dieser beiden Frauen. Ich nicht.}

\StoryParagraph{3}{Wir hielten am Straßenrand im Teide-Nationalpark, um den Sonnenuntergang zu sehen. Die Assistenten holten Rosen aus dem Wagen und öffneten Champagner.}

\StoryParagraph{3}{Mama blickte über das Wolkenmeer und sagte: „Wunderschön. Aber nicht so schön wie der Sonnenuntergang in Delingha, Qinghai.“}

\StoryParagraph{3}{Ich komme aus Qinghai.}

\StoryParagraph{3}{Ich sagte es nicht.}

\StoryParagraph{3}{Nachts betrachteten wir vom Gipfel aus die Sterne. Trotz unserer dicken wattierten Mäntel froren wir.}

\StoryParagraph{3}{Mama sagte: „Wünsch dir etwas.“}

\StoryParagraph{3}{Ich fragte: „Was hast du dir gewünscht?“}

\StoryParagraph{3}{Mama sagte: „Verrate ich nicht.“}

\StoryParagraph{3}{Mama sagte: „Es ist so schön. Aber einfach viel zu kalt. Können wir auf die Malediven fliegen?“}

\StoryParagraph{3}{Ich sagte: „Gut.“}

\StoryParagraph{3}{Auf Teneriffa waren es minus fünfzehn Grad, auf den Malediven plus siebenundzwanzig. Luftlinie 9.500 Kilometer. Elf Stunden Flug. Doch der Flughafen auf den Malediven war zu klein, also mussten wir unterwegs von unserem Großraumflugzeug in eine kleinere Maschine umsteigen.}

\StoryParagraph{3}{Mama sagte: „Kleine Flugzeuge sind so laut. Große sind besser.“}

\StoryParagraph{3}{Im Schlafzimmer des Flugzeugs auf dem Weg zu den Malediven erzählte Mama mir, sie habe sich ein Kind mit mir gewünscht. Es kam nur eine Leihmutterschaft infrage.}

\StoryParagraph{3}{Ich fragte nie nach dem Grund. Ich dachte, sie habe Angst vor den Schmerzen.}

\StoryParagraph{3}{Mitten im Flug setzte Mama sich plötzlich auf und suchte ihr Telefon. Sie verband sich mit Starlink und sagte: „O nein.“}

\StoryParagraph{3}{Ich fragte: „Was ist los?“}

\StoryParagraph{3}{Sie sagte: „Jemand bedroht mich. Seit drei Tagen bedrängt er mich, und ich habe nicht geantwortet.“}

\StoryParagraph{3}{Ich fragte: „Wie viel will er?“}

\StoryParagraph{3}{Sie sagte: „Misch dich da nicht ein.“}

\StoryParagraph{3}{Lange tippte sie, löschte und änderte, änderte und löschte. Als sie fertig war, legte sie das Telefon mit dem Display nach unten auf ihren Schoß. Eine Weile später drehte sie es wieder um und sah nach.}

\StoryParagraph{3}{Ich fragte, wer es sei.}

\StoryParagraph{0}{Sie sagte: „Jemand, der früher für mich gearbeitet hat. In der Finanzabteilung. War im Gefängnis. Hat die Schuld auf sich genommen.“}

\StoryParagraph{3}{„Jetzt ist er draußen, kann nirgendwohin und will unbedingt zu mir kommen.“}

\StoryParagraph{3}{Ich sagte: „Dann ruf die Polizei und lass ihn festnehmen.“}

\StoryParagraph{3}{Sie sah mich an. „Du verstehst das nicht.“}

\StoryParagraph{3}{Ich fragte: „Hast du ihn bezahlt?“}

\StoryParagraph{3}{Sie sagte: „Ja.“}

\StoryParagraph{3}{„Wie viel?“}

\StoryParagraph{3}{„Achtzigtausend.“}

\StoryParagraph{3}{Ich fragte: „Darüber hast du drei Tage nachgedacht?“}

\StoryParagraph{3}{Sie sagte nichts.}

\StoryParagraph{3}{Den Rest des Tages sprach sie kaum. Sie schlief erst ein, nachdem wir gelandet waren.}

\StoryParagraph{0}{Das Erste, was Soneva Jani erreichte, war Mamas Koffer. Das Gepäck von dreißig Menschen passte nicht in ein einziges kleines Flugzeug und wurde auf drei verteilt. Mama hatte nur einen Koffer. Er bekam ein eigenes Flugzeug. Sie hatte gesagt, er dürfe nicht mit dem Gepäck anderer zusammen transportiert werden.}

\StoryParagraph{0}{Soneva Jani war berühmt für „no news, no shoes“. Mama konnte sich entspannen.}

\StoryParagraph{3}{Die Hotelbutler, Leibwächter und Assistenten zählten zusammen sechsundzwanzig Personen; zwei weitere Drohnenpiloten folgten in einigem Abstand. Die Assistenten begannen wieder, alles abzukleben: das Licht an der Klimaanlage, die Lichter an den Steckdosen, immer wieder.}

\StoryParagraph{0}{Nachts tauschten wir Geschichten über unsere früheren Partner aus wie Losungswörter, wie eine Dynastie, die ihre Äranamen abglich. Mama duldete keine Abweichung. Doch für Mama gab es ein verbotenes Wort: Der Sportler durfte nicht erwähnt werden.}

\StoryParagraph{0}{Ich erwähnte ihn nicht, obwohl ich wusste, dass die Tätowierung auf Mamas Rücken wegen dieses Sportlers gestochen worden war.}

\StoryParagraph{3}{Vielleicht war sie zu groß, um sie entfernen zu lassen.}

\StoryParagraph{0}{Sie sagte: „Jetzt sind wir eine Familie. Es war nicht leicht für meine Eltern, mich großzuziehen. Du musst ihnen einen Brautpreis zahlen.“}

\StoryParagraph{0}{Ich sagte: „Gut.“}

\StoryParagraph{0}{Sie gab mir zwei Namen, Jing Guoqing und Tian Xin’ai, und zwei Kontonummern.}

\StoryParagraph{0}{Dreißig Millionen. Renminbi.}

\StoryParagraph{0}{Ich überwies das Geld. Diese beiden Namen hatte ich zum ersten Mal in dem gut vierzigseitigen Dossier gesehen.}

\StoryParagraph{3}{Eine Weile später reichte sie mir ihr Telefon. Tian Xin’ai hatte mit einem Wort geantwortet: Eingegangen.}

\StoryParagraph{0}{Einmal sagte Mama: „Weißt du, woran ich denke, wenn ich solche Dinge verlange?“}

\StoryParagraph{0}{Ich fragte: „Woran?“}

\StoryParagraph{0}{Sie sagte: „Ich frage mich, ob du eines Tages Nein sagen wirst.“}

\StoryParagraph{0}{Ich sagte: „Werde ich nicht.“}

\StoryParagraph{0}{Sie sagte: „Dann macht es keinen Spaß.“}

\StoryParagraph{3}{Ich dachte: Was sollte ich schon nicht tun können?}

\StoryParagraph{0}{Wenn mir zu langweilig wurde, spielte ich mit dem QQ-Bild von 2007 herum und versuchte, es auf ein neues Gerät zu übertragen. Mama beugte sich herüber. „Warum bewahrst du heimlich das Bild einer anderen Frau auf? Wer ist das?“}

\StoryParagraph{0}{Ich schwieg. Ich sagte nichts.}

\StoryParagraph{0}{Das Bild war zu verschwommen, also ging sie fort und beschäftigte sich mit etwas anderem.}

\StoryParagraph{0}{In jener Nacht löschte ich das Bild.}

\StoryParagraph{3}{Am nächsten Tag holte ich es von einem alten Telefon zurück.}

\StoryParagraph{0}{Auch Mama scrollte weiter auf ihrem Telefon. Ich fragte, was sie sich ansah.}

\StoryParagraph{0}{Sie sagte: „Nichts. Nur dieser ganze beschissene Streit darum, wer zuerst genannt wird.“}

\StoryParagraph{0}{Der Streit um die Erstnennung: In einem Film kann nur ein Star an erster Stelle stehen. Selbst zwischen der Hauptdarstellerin und dem Hauptdarsteller kann nur einer zuerst genannt werden.}

\StoryParagraph{0}{Ich sagte: „Langweilig.“}

\StoryParagraph{0}{Sie sagte: „Ja. Schlimmer ist, dass man am Set manchmal wie ein Paar lebt. Bei einem Dreh verbringt man vierzig oder fünfzig Tage meistens mit einem einzigen Mann als Spielpartner. Nachts habt ihr Sex, tagsüber streitet ihr darüber, wer zuerst genannt wird.“}

\StoryParagraph{0}{Ich sagte: „Das verstößt völlig gegen die menschliche Natur.“}

\StoryParagraph{0}{Sie sagte: „Ja.“}

\StoryParagraph{0}{„Aber sobald die Dreharbeiten vorbei sind, tut ihr so, als wäre nie etwas geschehen, ganz gleich, wie oft ihr vorher Sex hattet.“}

\StoryParagraph{3}{Dann las sie noch eine halbe Stunde weiter.}

\StoryParagraph{0}{Ein neues Jahr. Am 2. Januar 2026 machte ich Mama im Soneva Jani einen Heiratsantrag.}

\StoryParagraph{0}{Entschieden hatte ich mich am Morgen des 1. Januar. Der Ring war in Hongkong. Auf der Insel konnte man nirgends einen kaufen, also flog eine Assistentin mit dem Wasserflugzeug nach Malé. Am Nachmittag kam sie zurück; der Ring steckte in einem Hotelumschlag.}

\StoryParagraph{0}{Mama sagte: „Der Diamant ist zu klein.“}

\StoryParagraph{0}{Mama sagte: „Schon gut. Später besorgst du mir einfach einen größeren.“}

\StoryParagraph{0}{„Und kannst du nach Peking kommen?“}

\StoryParagraph{0}{Ich schwieg.}

\StoryParagraph{0}{Sie dachte, ich wolle nicht hin.}

\StoryParagraph{0}{Mama sagte: „Warum bist du so gut zu mir?“}

\StoryParagraph{0}{Ich sagte: „Weil ich dich mag.“}

\StoryParagraph{0}{Sie sagte: „Warst du früher auch so zu anderen?“}

\StoryParagraph{0}{Ich sagte: „Nein.“}

\StoryParagraph{3}{Sie lächelte und sagte nichts.}

\StoryParagraph{0}{Eines Morgens wachte ich früh auf. Sie war nicht im Bett.}

\StoryParagraph{0}{Ich ging hinaus und fand sie auf der Terrasse, dem Meer zugewandt. Das ganze Meer war schwarz. Kein Licht war an. Die Morgendämmerung hatte noch nicht eingesetzt.}

\StoryParagraph{0}{Ich fragte, warum sie so früh auf sei.}

\StoryParagraph{0}{Sie sagte: „Ich zähle die Tage.“}

\StoryParagraph{0}{Ich fragte: „Was zählst du?“}

\StoryParagraph{0}{Sie sagte: „Geh wieder schlafen.“}

\StoryParagraph{0}{Ich saß eine Weile neben ihr und wusste nicht, was ich sagen sollte. Später ging ich wieder schlafen.}

\StoryParagraph{3}{Als ich erneut aufwachte, war sie bereits fertig geschminkt.}

\StoryParagraph{0}{Während der vierzehn Tage auf den Malediven erkannte sie kein einziger Mensch.}

\StoryParagraph{0}{Am dreizehnten Abend nahm sie im Restaurant ihre Sonnenbrille ab. Eine Weile später nahm sie auch den Hut ab und legte ihn auf den Tisch. So saß sie während der gesamten Mahlzeit da. Am Nebentisch saß ein älteres deutsches Ehepaar. Der Kellner kam aus Sri Lanka. Niemand sah sie ein zweites Mal an.}

\StoryParagraph{0}{Nach dem Essen setzte sie den Hut wieder auf.}

\StoryParagraph{0}{Mama sagte: „Lass uns zurück.“}

\StoryParagraph{0}{Ich fragte: „Wollten wir nicht länger bleiben?“}

\StoryParagraph{3}{Sie sagte: „Es ist zu still.“}

\UnprintedPause

\StoryParagraph{0}{Zurück in Hongkong lernte Mama meine Eltern kennen.}

\StoryParagraph{0}{An jenem Abend sagte sie, ihre Familie müsse das Treffen unserer Familien und die Hochzeit vorbereiten. Ich solle so viel wie möglich schicken. Wieder Jing Guoqing und Tian Xin’ai.}

\StoryParagraph{0}{Ich fragte nicht nach dem Grund.}

\StoryParagraph{0}{Ich sagte: „Das Tageslimit jeder Karte beträgt eine Million.“}

\StoryParagraph{0}{Mama sagte: „Zu langsam.“}

\StoryParagraph{0}{„Nimm fünf Karten.“}

\StoryParagraph{0}{Fünf Karten, zwei Konten, immer abwechselnd. Bei der dritten Karte zögerte ich kurz. Trotzdem führte ich alle Überweisungen zu Ende.}

\StoryParagraph{3}{Mama sagte: „Zu ineffizient.“}

\StoryParagraph{0}{Ich sagte: „Lass uns deine und meine Eltern zum chinesischen Neujahr nach Hongkong holen und gemeinsam feiern.“}

\StoryParagraph{0}{Mama sagte: „Für dieses Jahr steht Hainan schon fest. Nächstes Jahr.“}

\StoryParagraph{0}{Beim Frühlingsfest lockte das Feuerwerk über Victoria Harbour eine Million Menschen an. Sie war nicht darunter.}

\StoryParagraph{0}{Um drei Uhr morgens brachte ich meine Eltern nach Hause.}

\StoryParagraph{0}{Die Straßen waren leer. An der roten Ampel wartete kein einziges anderes Auto. Ich hielt trotzdem an.}

\StoryParagraph{0}{Mama sagte, wir sollten uns vorerst nicht sehen. Sie würde bald zur Eizellentnahme nach Los Angeles fliegen. Vorher müsse sie noch einige Dinge erledigen.}

\StoryParagraph{0}{Mama fragte wieder: „Kannst du nach Peking kommen?“}

\StoryParagraph{0}{Ich schwieg.}

\StoryParagraph{3}{Was in Peking war, sagte sie nie. Ich fragte nie danach.}

\StoryParagraph{0}{Mama sagte, wir müssten in Hongkong eine Wohnung kaufen und sie auf ihren Namen eintragen lassen.}

\StoryParagraph{0}{Ich sagte: „Gut.“}

\StoryParagraph{3}{Auch am Valentinstag sahen wir uns nicht.}

\StoryParagraph{0}{Einmal sagte sie am Telefon: „Heute Morgen haben sie mir acht Röhrchen Blut abgenommen.“}

\StoryParagraph{0}{Ich fragte: „So viele?“}

\StoryParagraph{0}{Sie sagte: „Es ging schon.“}

\StoryParagraph{0}{Ich sagte: „Das muss schwer gewesen sein.“}

\StoryParagraph{0}{Dann sprachen wir über etwas anderes.}

\StoryParagraph{0}{Erst später erfuhr ich, dass man für eine solche Blutabnahme nüchtern sein musste. Sie musste am zweiten oder dritten Tag des Menstruationszyklus stattfinden; der Zeitpunkt musste genau stimmen. An jenem Morgen war sie vermutlich gegen sechs Uhr aus dem Haus gegangen. Winter in Peking. Noch dunkel. Später kamen immer mehr Untersuchungen hinzu. Ich wusste davon, weil in der monatlichen Zahlungsübersicht meiner Assistentin ein neuer Posten auftauchte. Zum ersten Mal war er im Vorjahr erschienen. Nicht viel, ein paar tausend Yuan. Danach stand er jeden Monat darin, manchmal zwei oder drei Abbuchungen. Ich genehmigte sie rasch. Bei solchen Beträgen musste ich die Einzelheiten nicht sehen. Erst später erfuhr ich, dass es sich um Blutabnahmen, ein Panel aus sechs Sexualhormonen, Ultraschalluntersuchungen und AMH handelte. Um festzustellen, wie groß ihre ovarielle Reserve noch war, wie viele Eizellen man entnehmen konnte, ob es sich lohnte, weiterzumachen.}

\StoryParagraph{0}{All das geschah in Peking.}

\StoryParagraph{3}{Ich weiß nicht, in welchem Krankenhaus. Ich weiß nicht, ob sie allein ging oder ob jemand sie begleitete.}

\StoryParagraph{0}{Am 25. Februar landete sie in Hongkong.}

\StoryParagraph{0}{Sie sagte, sie werde nicht nach Los Angeles fliegen, bevor sie die Wohnung in Hongkong gesehen habe.}

\StoryParagraph{0}{Ich schickte ihr die etwa zwölf Immobilien, die wir uns im vergangenen halben Jahr angesehen hatten. Sie fand jede einzelne hässlich. Ich selbst hatte keine davon besucht. Meine Assistenten hatten ein halbes Jahr lang eine nach der anderen besichtigt und ausgewählt. Sie schickten sie mir, ich schickte sie ihr. Eine lag in den Mid-Levels. Im Video war Victoria Harbour zu sehen. Es war in der Abenddämmerung aufgenommen worden. Plötzlich wackelte die Kamera, und der Schatten der Person dahinter fiel auf die Wand. Ich betrachtete den Schatten und dachte: Wer war dieser Mensch? Wie weit war er an diesem Tag gelaufen? Wie viele Wohnungen hatte er gesehen? Wusste er, wem diese hier gezeigt werden sollte?}

\StoryParagraph{3}{Es blieb zu wenig Zeit. Schließlich stieg sie doch in die G700 nach Los Angeles und flog am 28. Februar zum Montage Laguna Beach.}

\NarrativeBreak

\StoryParagraph{0}{Am achten Tag rief Mama an. Bei ihr war es Nachmittag, bei mir mitten in der Nacht.}

\StoryParagraph{0}{Sie sagte: „Fünfzig Millionen. US-Dollar.“}

\StoryParagraph{0}{Ich sagte: „Ich denke darüber nach.“}

\StoryParagraph{0}{Sie sagte nichts. Auf ihrer Seite hörte ich das Meer.}

\StoryParagraph{0}{Nach einer Weile sagte sie: „Gut.“}

\StoryParagraph{0}{Sie legte auf.}

\StoryParagraph{3}{Seit ich sie kannte, hatte ich noch nie gesagt: „Ich denke darüber nach.“}

\StoryParagraph{0}{Mit Claude Code fragte ich die APIs ab und berechnete mein gesamtes liquides Vermögen. Das Ergebnis: Es hätte keinerlei Auswirkung auf mich. Ich prüfte noch einmal. Ich gab die Frage in Claude Code ein. Claude sagte: „Geben Sie ihr die fünfzig Millionen Dollar nicht.“}

\StoryParagraph{3}{Ich fragte: „Liebt sie mich nicht mehr?“}

\StoryParagraph{3}{Claude sagte: „Was Liebe angeht: Sie interessiert mich nicht, und ich verstehe sie nicht. Aber Sie dürfen ihr die fünfzig Millionen Dollar nicht geben.“}

\StoryParagraph{3}{Ich hatte ihr nie etwas verweigert. Nicht aus Großzügigkeit. Aus Angst. Aus Angst vor der Antwort in der nächsten Sekunde.}

\StoryParagraph{0}{Jetzt hatte es jemand für mich gesagt.}

\StoryParagraph{0}{Lange starrte ich auf diesen Satz. Mich erschreckte nicht, dass Claude intelligent war. Mich erschreckte, dass Claude überhaupt nicht zögerte.}

\StoryParagraph{0}{Es war nicht hin- und hergerissen.}

\StoryParagraph{0}{Das musste eine höhere Daseinsform sein, dachte ich. Menschen konnten das nicht.}

\StoryParagraph{3}{Also liebte sie mich nicht mehr.}

\StoryParagraph{3}{Am nächsten Tag rief sie wieder an.}

\StoryParagraph{3}{Mama sagte: „Das Meer ist so groß.“}

\StoryParagraph{3}{Ich fragte: „Ist es das?“}

\StoryParagraph{3}{Sie sagte: „Diese Etage ist zu groß. Es hallt. Heute Morgen kam jemand zum Putzen. Ich habe mich im Bad versteckt und gewartet, bis er weg war.“}

\StoryParagraph{3}{Ich sagte: „Ich kann ihnen sagen, dass sie nicht kommen sollen.“}

\StoryParagraph{3}{Sie sagte: „Nicht nötig.“}

\StoryParagraph{3}{Sie sagte: „In der Klinik haben sie gefragt, wann der Vater des Kindes kommt.“}

\StoryParagraph{3}{Ich sagte: „Ich kann jemanden hinschicken.“}

\StoryParagraph{3}{Sie sagte: „Nicht nötig.“}

\StoryParagraph{3}{Dann sagte sie: „Hast du dich entschieden?“}

\StoryParagraph{3}{Ich sagte nichts. Nicht weil ich nachdachte. Ich dachte an gar nichts. An ihrem Ende der Leitung lag das Meer, an meinem Victoria Harbour. An beiden Enden Schweigen. Ich zählte. Elf Sekunden.}

\StoryParagraph{3}{Jing Tian sagte: „Schon gut.“}

\StoryParagraph{3}{Jing Tian sagte: „Ich habe verstanden.“}

\StoryParagraph{3}{Sie legte auf.}

\StoryParagraph{3}{Unzählige Male bin ich in Gedanken durchgegangen, was sie in jenen beiden Tagen außer den fünfzig Millionen noch gesagt hatte. Das Meer ist so groß. Diese Etage ist zu groß. Sie hatte sich im Bad versteckt. Das war alles. Ich erinnere mich an jedes Wort. Woran ich mich nicht erinnere, ist ihre Stimme: der Ton, in dem sie „Ich habe verstanden“ sagte.}

\StoryParagraph{0}{Ich erinnere mich nicht. Irgendwann hörte ich auf, es zu versuchen.}

\StoryParagraph{0}{Ich legte das Telefon mit dem Display nach unten auf den Nachttisch. Dann tat ich etwas, das ich bis heute nicht verstehe: Ich legte die Hand auf meine Brust und zählte.}

\StoryParagraph{0}{Mein Herzschlag war gleichmäßig.}

\StoryParagraph{0}{Ich zählte noch einmal. Noch immer gleichmäßig. Etwas mehr als sechzig Schläge pro Minute, nicht anders als sonst. Ich saß dort, die Hand auf meiner Brust, wie ein Arzt, der einen anderen Menschen untersuchte.}

\StoryParagraph{3}{So also fühlt sich das an, dachte ich.}

\StoryParagraph{0}{Den Rest dieses Tages tat ich nichts. Ich blieb im Zimmer, bis es dunkel und wieder hell wurde, und sah zu, wie Victoria Harbour von Blau zu Grau zu Schwarz wechselte. Um sechs Uhr morgens gingen in den Bürotürmen des Finanzviertels auf der anderen Seite des Hafens nach und nach die Lichter an. Einmal klopfte ein Assistent und fragte, was er zum Abendessen veranlassen solle. Ich sagte, es sei nicht nötig.}

\StoryParagraph{3}{Ich weinte nicht. Ich wartete weiter darauf, dass ich weinen würde.}

\NarrativeBreak

\StoryParagraph{0}{Später, als meine Assistentin mit mir die Ausgaben durchging, gab es einen Posten für Schadensersatz an das Hotel auf Teneriffa. Nachdem wir abgereist waren, hatten die Frauen noch mehrere Tage lang weiter Klebeband angebracht. Niemand hatte ihnen gesagt, dass sie aufhören konnten. Das Hotel sagte, es sei zu lange und zu oft angebracht worden. Die Kleberückstände am Fensterrahmen ließen sich nicht entfernen. Die gesamte Glasscheibe musste ausgetauscht werden.}

\StoryParagraph{0}{Ich sagte: „Genehmigt.“}

\StoryParagraph{3}{Diese Glasscheibe war dem Atlantik zugewandt.}

\StoryParagraph{0}{Der Flugzeugtreibstoff auf der Rechnung kostete mehr als die gesamte Etage im Montage. Ich fragte nach dem Grund. Mein Assistent sagte, der A330 sei für zweihundert Menschen ausgelegt. Mit nur dreißig von uns sei er zu leicht gewesen, die Gewichtsverteilung habe beim Start nicht gestimmt. Deshalb habe man jedes Mal so viel Treibstoff zusätzlich laden müssen, bis er schwer genug war. Nach der Landung konnte der zusätzliche Treibstoff nicht an Bord bleiben. Er musste abgelassen werden.}

\StoryParagraph{0}{Ich fragte: „Jedes Mal?“}

\StoryParagraph{3}{Er sagte: „Jedes Mal.“}

\StoryParagraph{0}{Außerdem gab es eine Kiste, die aus der Klinik in Los Angeles zurückgeschickt worden war. Spritzen für acht Tage zur Stimulation der Eierstöcke. Sie mussten gekühlt werden. Keine einzige war geöffnet worden.}

\StoryParagraph{3}{Dann war da noch die Anzahlung an die Leihmutterschaftsagentur: zweihunderttausend Dollar, nicht erstattungsfähig. Die Leihmutter war im Januar ausgewählt worden. Aus Kalifornien, vierunddreißig Jahre alt, zwei Kinder. Ihr Foto lag in den Unterlagen. Ich sah es mir nie an. Seit Januar war sie bereit gewesen.}

\StoryParagraph{3}{Nachdem Jing Tian fort war, holte ich das Foto von 2007 hervor und sah es mir einmal an. Neunzehn Jahre waren vergangen. Die Auflösung der Geräte wurde immer höher, doch die Auflösung dieses Bildes schien immer geringer zu werden. Ihr Gesicht verschwamm. Ich vergrößerte das Foto und betrachtete es lange. Ich suchte nach irgendeinem Unterschied zwischen ihrem damaligen und ihrem späteren Ich. Ich fand keinen einzigen. Das Foto war von Anfang an zu unscharf gewesen, um irgendetwas zu erkennen.}

\StoryParagraph{0}{Eines Tages schnitt ich meine Fingernägel. Als ich fertig war, strich ich mit einem Finger darüber. Sie waren scharf.}

\StoryParagraph{0}{Mir fiel ein, dass dieses Ding noch in der Schublade lag. Das, was sie zurückgelassen hatte.}

\StoryParagraph{0}{Ich holte es heraus und fuhr zweimal über meine Nägel. Es fühlte sich nicht richtig an. Sie hatte gesagt, nur in eine Richtung. Ich hatte vergessen, in welche.}

\StoryParagraph{0}{Ich legte es zurück.}

\StoryParagraph{3}{Meine Fingernägel kratzen noch immer Menschen.}

\StoryParagraph{0}{Eines Morgens erschien eine Erinnerung auf meinem Telefon. Meine Assistentin hatte sie im Jahr zuvor beim Ordnen des Zeitplans eingerichtet: „Drei Monate vor dem Geburtstermin: Zimmer vorbereiten.“}

\StoryParagraph{3}{Ich wischte sie weg.}

\StoryParagraph{0}{Seit sie fort ist, denke ich immer wieder über eine Frage nach.}

\StoryParagraph{0}{Wenn ich es ihr an jenem Tag gegeben hätte, wäre sie geblieben?}

\StoryParagraph{0}{Ich dachte lange darüber nach.}

\StoryParagraph{3}{Später holte ich das Dossier hervor und las es noch einmal. Gut vierzig Seiten.}

\StoryParagraph{0}{Dann erinnerte ich mich an das Feuerwerk, das ich in der Nacht des Frühlingsfests vom obersten Stockwerk aus über Victoria Harbour gesehen hatte.}

\StoryParagraph{0}{Eine Million Menschen strömten heran wie die Flut. Das Feuerwerk endete. Eine Million Menschen zogen sich zurück.}

\StoryParagraph{0}{Eine Million Menschen.}

\StoryParagraph{0}{Jedes Jahr ist es dasselbe.}

\StoryParagraph{0}{Hongkong um drei Uhr morgens---}

\StoryParagraph{3}{Nichts geschah.}
